\section{Phương trình, bất phương trình mũ và logarit}

\subsection{Đề 1}
\Opensolutionfile{ans}[ans/ans-1-C6-B4]
\caulc

\begin{ex}%[1D6H4-3]
	Tập nghiệm của bất phương trình $\log_3 \left(x+1\right) < 2$ là
	\choice
	{\True $\left(-1; 8\right)$}
	{$\left(-\infty; 8\right)$}
	{$\left(-\infty; 7\right)$}
	{$\left(-1; 7\right)$}
	\loigiai{Điều kiện $x>-1$.\\
		Khi đó bất phương trình tương đương
		$\log_3 \left(x+1\right) < 2\Leftrightarrow x+1<3^2\Leftrightarrow x<8$.\\
		Vậy $-1<x<8$.	}
\end{ex}
%%%==============EX_2============%%%
\begin{ex}%[1D6H4-4]
	Nghiệm của phương trình $\log_3 \left(2x-1\right)=2$ là
	\choice
	{$x=3$}
	{\True $x=5$}
	{$x=\dfrac{9}{2}$}
	{$x=\dfrac{7}{2}$}
	\loigiai{Điều kiện $x>\dfrac{1}{2}$.\\
		Khi đó bất phương trình tương đương
		$\log_3 \left(2x-1\right)=2\Leftrightarrow 2x-1=3^2\Leftrightarrow x=5$.\\
		Vậy $-x=5$.
	}
\end{ex}
\begin{ex}%[1D6H4-4]
	Tập nghiệm $S$ của phương trình $3^{x^2-2x}=27$.
	\choice
	{$S=\left\{1; 3\right\}$}
	{$S=\left\{-3; 1\right\}$}
	{\True $S=\left\{-1; 3\right\}$}
	{$S=\left\{-3;-1\right\}$}
	\loigiai{
		$3^{x^2-2x}=27\Leftrightarrow x^2-2x=3\Leftrightarrow \hoac{&x=-1\\&x=3.}$\\
		Vậy tập nghiệm phương trình $S=\left\{-1; 3\right\}$.	}
\end{ex}
\begin{ex}%[1D6H4-2]
	Gọi $x_1, \; x_2$ là hai nghiệm của phương trình $7^{x+1}=\left(\dfrac{1}{7} \right)^{x^2-2x-3}$. Khi đó $x_1^2+x_2^2$ bằng
	\choice
	{$17$}
	{\True $5$}
	{$3$}
	{$1$}
	\loigiai{
		$7^{x+1}=\left(\dfrac{1}{7} \right)^{x^2-2x-3}\Leftrightarrow-x-1=x^2-2x-3\Leftrightarrow x^2-x-2\Leftrightarrow \hoac{&x=-1\\&x=2.}$\\
		Khi đó 	$x_1^2+x_2^2=5$.}
\end{ex}
%%%==============EX_8============%%%
\begin{ex}%[1D6H4-4]
	Số nghiệm của phương trình $\log_3 \left(6+x\right)+\log_3 9x-5=0$ là 
	\choice
	{$2$}
	{\True $1$}
	{$0$}
	{$3$}
	\loigiai{Điều kiện xác định của phương trình $x>0$.
		\begin{eqnarray*}
			& & \log_3 \left(6+x\right)+\log_3 9x-5=0\\
			&\Leftrightarrow & \log_3\left(9x(6+x) \right)=5\\
			&\Leftrightarrow & x(6+x)=27\Leftrightarrow x^2+6x-27=0\Leftrightarrow\hoac{&x=-9\\&x=3.} 
		\end{eqnarray*}
		Vậy so với điều kiện ta chọn nghiệm $x=3$.}
\end{ex}
\begin{ex}%[1D6H4-5]
	Biết tập nghiệm của bất phương trình $2^x < 3-\dfrac{2}{2^x}$ là $\left(a; b\right)$. Giá trị $5a+3b$ bằng
	\choice
	{$2$}
	{$1$}
	{\True $3$}
	{$0$}
	\loigiai{Đặt $t=2^2$, $t>0$. Khi đó bất phương trình viết lại $$t<3-\dfrac{2}{t}\Leftrightarrow t^2-3t+2<0\Leftrightarrow 1<t<2\Leftrightarrow 0<x<1.$$
		Khi đó $a=0$, $b=1$. Vậy $5a+3b=3$.
	}
\end{ex}
\begin{ex}%[1D6H4-4]
	Tìm tất cả các giá trị của tham số $m$ để phương trình $2^{2x+1}-2^{x+3}-2m=0$ có hai nghiệm phân biệt.	
	\choice
	{$m > 0$}
	{$m >-4$}
	{\True $-4< m < 0$}
	{$m <-4$}
	\loigiai{
		Đặt $t=2^x, \left(t > 0\right)$.\\
		Phương trình trở thành $2t^2-8t-2m=0$ (*)\\
		Phương trình có 2 nghiệm phân biệt khi phương trình (*) có $2$ nghiệm dương phân biệt $$\Leftrightarrow \heva{&{\Delta ' > 0} \\
			&{S> 0} \\
			&{P> 0}} \Leftrightarrow \heva{&{16+4m > 0} \\
			&{m < 0}
			&} \Leftrightarrow-4< m < 0.$$
	}
\end{ex}

\begin{ex}%[1D6H4-4]
	Phương trình $9^x-3\cdot 3^x+2=0$ có 2 nghiệm $x_1$, $x_2$ với $x_1 < x_2$. Giá trị $2x_1+3x_2$ bằng
	\choice
	{$4\log_2 3$}
	{$2$}
	{$0$}
	{\True $3\log_3 2$}
	\loigiai{$9^x-3\cdot 3^x+2=0\Leftrightarrow \hoac{&3^x=1\\&3^x=2}\Leftrightarrow \hoac{&x=0\\&x=\log_32.}$\\
		Khi đó $2x_1+3x_2=3\log_3 2.$	}
\end{ex}

\begin{ex}%[1D6H4-5]
	Tập nghiệm $S$ của bất phương trình $5^{x+2} < \left(\dfrac{1}{25} \right)^{-x}$ là
	\choice
	{$S=\left(-\infty;2\right)$}
	{$S=\left(-\infty;1\right)$}
	{$S=\left(1;+\infty \right)$}
	{\True $S=\left(2;+\infty \right)$}
	\loigiai{Tập xác định $\mathscr{D}=\mathbb{R}$.\\
		
	}
\end{ex}
\begin{ex}%[1D6H4-6]
	Cho biết chu kì bán rã của chất phỏng xạ plutonium $Pu^{239}$ là $24.360$ năm (tức là lượng $Pu^{239}$ sau $24360$ năm phân hủy thì chỉ còn lại một nửa
	Sự phân hủy được tính bởi công thức $S=Ae^{rt}$, trong đó $A$ là lượng chất phóng xạ ban đầu, $r$ là tỉ lệ phân hủy hàng năm ($r < 0$), $t$ là thời gian phân huỷ, $S$ là lượng còn lại sau thời gian phân hủy $t$. Sau khoảng bao nhiêu năm thì 10 gam $Pu^{239}$ phân hủy còn lại $1$ gam?
	\choice
	{$833335$}
	{\True $822235$}
	{$800235$}
	{$822000$}
	\loigiai{
		Vì chu kì bán rã của chất phỏng xạ plutonium $Pu^{239}$ là $24.360$ năm nên ta có:
		\[5=10e^{r\cdot24360} \Leftrightarrow r\approx-0,000028.
		\]
		Sự phân hủy của chất phỏng xạ plutonium $Pu^{239}$ được tính bởi công thức $S=Ae^{-0,000028t}$.\\	
		Với $S=1$ thì $1=10e^{-0,000028t}$ $\Leftrightarrow t\approx 822235$ năm.\\
		Vậy sau khoảng $822235$ năm thì $10$ gam $Pu^{239}$ phân hủy còn lại $1$ gam.
	}
\end{ex}
\begin{ex}%[1D6V4-5]
	Xét bất phương trình $\log_2^2 2x-2\left(m+1\right)\log_2 x-2< 0$. Tìm tất cả các giá trị của tham số $m$ để bất phương trình có nghiệm thuộc khoảng $\left(\sqrt{2;}+\infty \right)$.
	\choice
	{$m\in \left(0;+\infty \right)$}
	{$m\in \left(-\dfrac{3}{4};0\right)$}
	{\True $m\in \left(-\dfrac{3}{4};+\infty \right)$}
	{$m\in \left(-\infty;0\right)$}
	\loigiai{
		Điều kiện: $x > 0$;
		\[\log_2^2 2x-2\left(m+1\right)\log_2 x-2< 0 \Leftrightarrow \left(1+\log_2 x\right)^2-2\left(m+1\right)\log_2 x-2 < 0 \; \; \; (1).
		\]
		Đặt $t=\log_2 x$.Vì $x > \sqrt{2}$ nên $\log_2 x > \log_2 \sqrt{2}=\dfrac{1}{2}$. Do đó $t\in \left(\dfrac{1}{2};+\infty \right)$.\\
		$(1)$ thành $\left(1+t\right)^2-2\left(m+1\right)t-2< 0$ $\Leftrightarrow t^2-2mt-1< 0$ $(2)$.\\
		Yêu cầu bài toán tương đương tìm $m$ để bpt (2) có nghiệm thuộc $\left(\dfrac{1}{2};+\infty \right)$.\\
		Xét bất phương trình (2) có: $\Delta '=m^2+1> 0, \; \forall m\in \mathbb{R}$.\\
		$f(t)=t^2-2mt-1=0$ có $ac < 0$ nên (2) luôn có 2 nghiệm phân biệt $t_1 < 0< t_2$.\\
		Khi đó cần $\dfrac{1}{2} < t_2 \Leftrightarrow m+\sqrt{m^2+1} > \dfrac{1}{2} \Leftrightarrow m >-\dfrac{3}{4}$.
	}
\end{ex}
%%%==============EX_38============%%%
\begin{ex}%[1D6V4-5]
	Cho bất phương trình $1+\log_5 \left(x^2+1\right)\ge \log_5 \left(mx^2+4x+m\right) (1)$. Tất cả các giá trị của $m$ để $(1)$ nghiệm đúng với mọi số thực $x$ là
	\choice
	{$2\le m\le 3$}
	{\True $2< m\le 3$}
	{$-3\le m\le 7$}
	{$m\le 3$; $m\ge 7$}
	\loigiai{
		Yêu cầu bài toán tương đương với
		\[\heva{&{mx^2+4x+m > 0\; \; \; \; \; \; \; \; \; \forall x\in \mathbb{R} \; \; \;} \\
			&{\log_5\left(5x^2+5\right)\ge \log_5\left(mx^2+4x+m\right) \; \; \forall x\in \mathbb{R} \; \;}}
		\]
		\[\Leftrightarrow \heva{&{mx^2+4x+m > 0\; \; \; \; \; \; \; \; \; \forall x\in \mathbb{R} \; \; \; (1)} \\
			&{\left(5-m\right)x^2-4x+\left(5-m\right)\ge 0\; \; \forall x\in \mathbb{R} \; \; (2)}}
		\]
		Với $(1)\Leftrightarrow \heva{&{m > 0} \\
			&{4-m^2< 0}
			&} \Leftrightarrow m > 2$.\\
		Xét $(2)$, dễ thấy $m=5$ mệnh đề không đúng. Khi $m\ne 5$
		\begin{eqnarray*}
			(2)&\Leftrightarrow  & \heva{&{5-m > 0} \\
				&{4-\left(5-m\right)^2\le 0}
				&}\\
			&\Leftrightarrow & \heva{&{m < 5} \\
				&{-m^2+10m-21\le 0}
				&} \\
			&\Leftrightarrow & \heva{&{m < 5} \\
				&{m\in (-\infty;3]\cup [7;+\infty)}
				&} \Leftrightarrow m\in (-\infty;3].
		\end{eqnarray*}		
		Vậy $2< m\le 3$.
	}
\end{ex}









\Closesolutionfile{ans}
\indapan{6}{ans/ans-1-C6-B4}
\cauds
\Opensolutionfile{ans}[ans/ans-1-C6-B4-DS]


\begin{ex}%[1D6H4-2]
	Giải được các phương trình sau. Khi đó:
	\choiceTF{Phương trình $3^{x-1}=9$ có một nghiệm}
	{Phương trình $5^{x-1}=\left(\dfrac{1}{25} \right)^x$ có nghiệm lớn hơn 3}
	{Phương trình $3^{x-2}=6$ có chung tập nghiệm với phương trình $x^2-2x+4=0$}
	{Phương trình $7^{x+2}-40.7^x=9$ có một nghiệm $x=a$, khi đó ${\mathop{\lim}\limits_{x\to a}} \left(x^2+2x+5\right)=6$}
	\loigiai{
		\begin{itemchoice}
			\itemch $3^{x-1}=9\Leftrightarrow 3^{x-1}=3^2 \Leftrightarrow x-1=2\Leftrightarrow x=3$.\\
			Vậy phương trình có nghiệm là $x=3$.
			\itemch $5^{x-1}=\left(\dfrac{1}{25} \right)^x \Leftrightarrow 5^{x-1}=5^{-2x} \Leftrightarrow x-1=-2x\Leftrightarrow x=\dfrac{1}{3}$.\\
			Vậy phương trình có nghiệm là $x=\dfrac{1}{3}$.
			\itemch $3^{x-2}=6\Leftrightarrow x-2=\log_3 6\Leftrightarrow x=\log_3 6+2$.
			\\	Vậy phương trình có nghiệm là $x=\log_3 6+2$.
			\itemch $7^{x+2}-40.7^x=9\Leftrightarrow 7^2.7^x-40.7^x=9\Leftrightarrow 9.7^x=9\Leftrightarrow 7^x=1\Leftrightarrow x=0$.\\Vậy phương trình có nghiệm là $x=0$.\\
			Suy ra ${\mathop{\lim}\limits_{x\to 0}} \left(x^2+2x+5\right)=5$.
		\end{itemchoice}
	}
\end{ex}

\begin{ex}%[1D6H4-2]
	Cho phương trình $\log_3 (x+6)=\log_3 (x-1)+1$ (*). Khi đó
	\choiceTF{\True Điều kiện $x > 1$}
	{Phương trình (*) có chung tập nghiệm với phương trình $\dfrac{x^2-11x+9}{x-1}=0$}
	{Gọi $x=a$ là nghiệm của phương trình (*), khi đó ${\mathop{\lim}\limits_{x\to a}} \left(x-3\right)=\dfrac{5}{2}$}
	{Nghiệm của phương trình (*) là hoành độ giao điểm của đường thẳng $\break d_1\colon2x-y-8=0$ với $\break d_2\colon y=0$}
	\loigiai{
		Điều kiện $\heva{&{x+6> 0} \\
			&{x-1> 0}}\Leftrightarrow x > 1$.
		\[\log_3 (x+6)=\log_3 (x-1)+1\Leftrightarrow \log_3 (x+6)=\log_3 (x-1)+\log_3 3
		\]
		$\Leftrightarrow \log_3 (x+6)=\log_3 3(x-1)\Rightarrow x+6=3(x-1)\Leftrightarrow x=\dfrac{9}{2}$ (thoả mãn điều kiện).\\	
		Vậy phương trình có nghiệm là $x=\dfrac{9}{2}$.
		\begin{itemchoice}
			\itemch Đúng.
			\itemch Sai. Vì $\dfrac{x^2-11x+9}{x-1}=0$ không có nghiệm $x=\dfrac{9}{2}$.
			\itemch Sai. Vì  ${\mathop{\lim}\limits_{x\to a}} \left(x-3\right)=\dfrac{-4}{2}=-2$.
			\itemch Sai. Vì  $x=4$ là hoành độ giao điểm của đường thẳng $d_1\colon2x-y-8=0$ với $d_2\colon y=0$.
		\end{itemchoice}	
	}
\end{ex}
%%%==============EX_7============%%%
\begin{ex}%[1D6H4-3]
	Giải được các bất phương trình sau. Khi đó
	\choiceTF{$16^x < \dfrac{1}{4}$ có tập nghiệm là $\left(-\infty;-\dfrac{1}{2} \right)$}
	{$5^{x-1} \ge \left(\dfrac{1}{25} \right)^x$ có nghiệm lớn nhất là $x=\dfrac{1}{3}$}
	{$(0, 3)^{x-2} \le 3$ có nghiệm lớn nhất là $x=2+\log_6 3$}
	{$2.7^{x+2} > 9$ có tập nghiệm là $\left(-2+\log_7 \left(\dfrac{9}{2} \right);+\infty \right)$}
	\loigiai{
		\begin{itemchoice}
			\itemch $16^x < \dfrac{1}{4} \Leftrightarrow 2^{4x} < 2^{-2} \Leftrightarrow 4x <-2\Leftrightarrow x <-\dfrac{1}{2}$ (do $2 > 1$).
			
			Vậy nghiệm của bất phương trình là $x <-\dfrac{1}{2}$
			\itemch $5^{x-1} \ge \left(\dfrac{1}{25} \right)^x \Leftrightarrow 5^{x-1} \ge 5^{-2x} \Leftrightarrow x-1\ge-2x\Leftrightarrow x\ge \dfrac{1}{3} ($ do $5 < 1)$.
			
			Vậy nghiệm của bất phương trình là $x\ge \dfrac{1}{3}$
			\itemch $(0, 3)^{x-2} \le 3\Leftrightarrow x-2\ge \log _{0, 3} 3\Leftrightarrow x\ge 2+\log _{0, 3} 3$ (do $0 < 0, 3 < 1$).
			
			Vậy nghiệm của bất phương trình là $x\ge 2+\log _{0, 3} 3$
			\itemch $2.7^{x+2} > 9\Leftrightarrow 7^{x+2} > \dfrac{9}{2} \Leftrightarrow x+2 > \log_7 \left(\dfrac{9}{2} \right)\Leftrightarrow x >-2+\log_7 \left(\dfrac{9}{2} \right)(\text{vì} \, 7 > 1)$.\\
			Vậy nghiệm của bất phương trình là $x >-2+\log_7 \left(\dfrac{9}{2} \right)$.
		\end{itemchoice}			
	}
\end{ex}

\begin{ex}%[1D6V4-5]
	Cho bất phương trình $\left(\dfrac{1}{9} \right)^x \ge 27\cdot3^x$, có tập nghiệm là $S=\left(a;b\right]$. Khi đó
	\choiceTF{ \True Bất phương trình có chung tập nghiệm với $3^{-2x} \ge 3^{3+x}$}
	{\True Có $A\left(0;b\right)$ giao điểm của đồ thị $y=x^3+2x-1$ với trục tung $Oy$}
	{\True ${\mathop{\lim}\limits_{x\to a}} \left(3x+2\right)=a$}
	{${\mathop{\lim}\limits_{x\to b}} \left(3x+2\right)=2$}
	\loigiai{
		\begin{eqnarray*}
			& & \left(\dfrac{1}{9} \right)^x \ge 27\cdot3^x \\
			&\Leftrightarrow & 3^{-2x} \ge 3^3 \cdot 3^x \\
			&\Leftrightarrow & 3^{-2x} \ge 3^{3+x}\Leftrightarrow-2x\ge 3+x (\text{vì}\, 3 > 1) \Leftrightarrow x\le-1.
		\end{eqnarray*}	
		Vậy nghiệm của bất phương trình là $x\le-1$.
		\begin{itemchoice}
			\itemch Đúng.
			\itemch Theo trên $a$ là vô cùng, $b=-1$.\\
			Khi đó $A\left(0;-1\right)$ giao điểm của đồ thị $y=x^3+2x-1$ với trục tung $Oy$.
			\itemch Đúng. Vì ${\mathop{\lim}\limits_{x\to -\infty}} \left(3x+2\right)=-\infty$.
			\itemch Sai ${\mathop{\lim}\limits_{x\to -1}} \left(3x+2\right)=-1$.
	\end{itemchoice}		}
\end{ex}





\Closesolutionfile{ans}
\indapan{3}{ans/ans-1-C6-B4-DS}

\Opensolutionfile{ans}[ans/ans-1-C6-B4-KQ]
\caukq

\begin{ex}%[1D6V4-6]
	Dân số ở một địa phương được ước tính theo công thức $S=A\cdot e^{r.t}$, trong đó $A$ không đổi là dân số của năm 2023, $S$ là dân số sau $t$ năm, $r$ là tỉ lệ tăng dân số hằng năm. Hỏi đến năm nào thì dân số ở địa phương đó sẽ đạt gấp đôi dân số năm 2023? Biết $r=1,13\% /$ năm.
	\shortans{$2085$}	
	\loigiai{
		Dân số đạt gấp đôi nghĩa là $S=2A$, ta có
		\[2A=A\cdot e^{1,13\%.t} \Leftrightarrow e^{1,13\%.t}=2\Leftrightarrow 1,13\%.t=\ln_e 2\Leftrightarrow t=\dfrac{\ln 2}{1,13\%} \approx 61,34 \; (do\; e > 1 \;).\;
		\]
		Vậy sau $62$ năm tức đến năm $2085$ thì dân số ở địa phương đó sẽ gấp đôi dân số năm $2023$.
	}
\end{ex}

\begin{ex}%[1D6H4-3]
	Tìm nghiệm phương trình $\log _{\dfrac{1}{4}} (-x+2)=-2$
	\shortans{$-14$}	
	\loigiai{
		Điều kiện: $-x+2> 0\Leftrightarrow x < 2$.
		
		$\log _{\dfrac{1}{4}} (-x+2)=-2\Leftrightarrow-x+2=\left(\dfrac{1}{4} \right)^{-2} \Leftrightarrow x=-14$ (thoả mãn điều kiện).
		
		Vậy phương trình có nghiệm là $x=-14$.
	}
\end{ex}
\begin{ex}%[1D6H4-4]
	Tìm nghiệm của phương trình $\ln (x+1)+\ln (x+3)=\ln (x+7)$.
	\shortans{$1$}	
	\loigiai{
		Với điều kiện $x >-1$, phương trình trở thành:
		\[\ln (x+1)(x+3)=\ln (x+7)\Leftrightarrow (x+1)(x+3)=x+7\Leftrightarrow \left[\begin{array}{l} {x=1} \\ {x=-4} \end{array}\right..
		\]
		Kết hợp điều kiện ta có nghiệm của phương trình là $x=1$.
	}
\end{ex}
\begin{ex}%[1D6V4-3]
	Tập nghiệm bất phương trình $\log_2 3+\log_5 x\ge 1+\log_2 3\cdot \log_5 x$ là $(a;b)$. Khi đó $b-a$ là 
	\shortans{$5$}	
	
	\loigiai{
		Điều kiện: $x > 0.(*)$
		
		Khi đó, do $1-\log_2 3< 0$ và cơ số $5> 1$ nên bất phương trình trở thành:
		\[\log_5 x\le 1\Rightarrow x\le 5.\;
		\]
		Kết hợp với điều kiện $(*)$, ta được nghiệm của bất phương trình là $0< x\le 5$.
	}
\end{ex}

\begin{ex}%[1D6V4-4]
	Cho phương trình $\left(2\log_3^{2}x-\log_3x-1 \right)\sqrt{3^x-m}=0$ ($m$ là tham số thực). Có tất cả bao nhiêu giá trị nguyên dương của $m$ để phương trình đã cho có đúng hai nghiệm phân biệt.
	\shortans{$26$}	
	\loigiai
	{Điều kiện $\heva{&x>0\\&3^x-m\geq 0}$. Khi đó phương trình tương đương $\hoac{&\log_3{x}=1\\&\log_3{x}=-\dfrac{1}{2}\\&3^x=m} \Leftrightarrow \hoac{&x=3\\&x=\dfrac{1}{\sqrt{3}}\\&3^x=m.}$\\
		Vì $m$ là số nguyên dương nên ta xét các trường hợp sau:
		\begin{itemize}
			\item Trường hợp $1$: Với $m=1$, phương trình có hai nghiệm phân biệt $x=3$, $x=\dfrac{1}{\sqrt{3}}$.
			\item Trường hợp $2$: Với $m>1$, khi đó phương trình $3^x=m$ luôn có nghiệm dương $x=\log_3{m}$. Do đó phương trình đã cho có đúng hai nghiệm phận biệt khi và chỉ khi
			$$\dfrac{1}{\sqrt{3}}\leq \log_3{m}<3 \Leftrightarrow 3^{\tfrac{1}{\sqrt{3}}} \leq m <27.$$
		\end{itemize}
		Do $m$ nguyên dương nên phương trình đã cho có đúng hai nghiệm phân biệt khi và chỉ khi $m\in \{1;2;\ldots;26\}$. Vậy có $26$ giá trị của $m$ thỏa mãn yêu cầu bài toán.
	}
\end{ex}
\begin{ex}%[1D6C4-5]
	Cho $x\geq 0, y\geq 0, x+y>0$ thỏa mãn $2^{x^2+y^2}+2023^{x+y} \cdot \log _2 \dfrac{x^2+y^2}{x+y} \leq 4^{x+y}+2023^{x+y}$. Tìm tổng giá trị lớn nhất và giá trị nhỏ nhất của biểu thức $P=x^2+y^2-8x-2y+10$ (làm tròn đến hàng phần trăm)
	\shortans{$5{,}51$}		
	\loigiai{
		$$
		\begin{aligned}
			& 2^{x^2+y^2}+2023^{x+y} \cdot \log _2 \dfrac{x^2+y^2}{x+y} \leqslant 4^{x+y}+2023^{x+y}  \\
			& \Leftrightarrow 2^{x^2+y^2}+2023^{x+y} \cdot \log _2 \dfrac{x^2+y^2}{2(x+y)} \leq 4^{x+y}  \\
			& \Leftrightarrow 2023^{x+y} \cdot\left[\log _2\left(x^2+y^2\right)-\log _2[2(x+y)]\right] \leq 2^{2(x+y)}-2^{x^2+y^2}
		\end{aligned}
		$$
		Nếu $x^2+y^2>2(x+y)$. Khi đó $VT>0 ; VP<0$ (không thoả mãn).\\
		Nếu $x^2+y^2 \leq 2(x+y)$. Khi đó $VT \leq 0 ; VP \geq 0$ luôn thoả mãn.\\
		Vậy $x^2+y^2-2 x-2 y \leq 0 \Rightarrow(x;y)$ thuộc phần hình tròn tâm $I(1;1)$ bán kính $r=\sqrt{2}$ (với $x\geq0, y\geq 0, x+y>0$).\\
		Khi đó
		$P=x^2+y^2-8x-2y+10 \Rightarrow(x;y)$ thuộc phần đường tròn tâm $K(4;1)$ bán kính $R=\sqrt{P+7}$ thoả mãn $x\geq 0, y\geq 0, x+y>0; \mathrm{d}=KI=3$.
		\begin{center}
			\begin{tikzpicture}[x=1cm,y=1cm,scale=1]
				\draw[->] (-1,0)--(7,0)node[below right]{$x$};
				\draw[->] (0,-2)--(0,4)node[left]{$y$};
				\fill (0,0)node[below left]{ $O$};
				\path (1,0)coordinate[label=below:$1$](A) (2,0)coordinate[label=below:$2$](B) (3,0)coordinate[label=below:$3$](C) (4,0)coordinate[label=below:$4$](D) (5,0)coordinate[label=below:$5$](E) (6,0)coordinate[label=below:$6$](F) (0,-1)coordinate[label=left:$-1$](G) (0,1)coordinate[label=left:$1$](H) (0,2)coordinate[label=left:$2$](N) (0,3)coordinate[label=left:$3$](M) (1,1)coordinate[label=above:$I$](I) (4,1)coordinate[label=above:$K$](K) (1.8,0)coordinate(T);
				\draw[black] let \p1=($(N)-(I)$) in (I) circle ({veclen(\x1,\y1)});
				\draw[black] let \p1=($(T)-(K)$) in (K) circle ({veclen(\x1,\y1)});
				\coordinate (W) at ($(I)!1!45:(B)$);
				
				\draw[dashed] (W)--(K)--(N);
				
				\foreach \diem in {A,B,C,D,E,F,G,H,M,N,I,K,W} \fill[black](\diem)circle(1.5pt);
				
			\end{tikzpicture}
		\end{center}		
		Dựa vào hình vẽ, để tồn tại $(x;y)$ ta phải có $\mathrm{d}-r\leq R\leq KA,(A(0;2)) \Leftrightarrow 3-\sqrt{2} \leq \sqrt{P+7} \leq \sqrt{17}$
		$$
		\begin{aligned}
			& 11-6\sqrt 2\leq P+7\leq 17\Leftrightarrow 4-6\sqrt 2\leq P\leq 10\\ &	 P _{\max}+P _{\min}=10+4-6\sqrt 2=14-6\sqrt 2\approx5{,}51.
		\end{aligned}
		$$		
	}
\end{ex}

\Closesolutionfile{ans}
\indapan{6}{ans/ans-1-C6-B4-KQ}











\subsection{Đề 2}
\Opensolutionfile{ans}[ans/ans-1-C6-B4]
\caulc

\begin{ex}%[1D6H4-2]
	Tính tổng tất cả các nghiệm của phương trình $2^{2x^2+5x+4}=4$.
	\choice
	{$-1$}
	{$1$}
	{$\dfrac{5}{2}$}
	{\True $-\dfrac{5}{2}$}
	\loigiai{
		$2^{2x^2+5x+4}=4\Leftrightarrow 2x^2+5x+4=2\Leftrightarrow2x^2+5x+2=0\Leftrightarrow\hoac{&x=-2\\&x=-\dfrac{1}{2}.}$\\
		Khi đó tổng các nghiệm của phương trình trên là $-\dfrac{5}{2}$.	}
\end{ex}

\begin{ex}%[1D6H4-4]
	Tổng các nghiệm của phương trình $\log _{\sqrt{3}} \left(x-2\right)+\log_3 \left(x-4\right)^2=0$ là $S=a+b\sqrt{2}$ (với $a, b$ là các số nguyên). Giá trị của biểu thức $Q=a+3b$ bằng
	\choice
	{$7$}
	{\True $9$}
	{$6$}
	{$8$}
	\loigiai{Điều kiện xác định $\heva{&x>2\\&x\not=4.}$\\
		Phương trình đã cho tương đương với
		\begin{eqnarray*}
			& & \log_3(x-2)^2+\log_3(x-4)^2=0\\
			&\Leftrightarrow & \log_3(x-2)^2(x-4)^2=0\\
			&\Leftrightarrow & (x-2)^2(x-4)^2=1\Leftrightarrow (x-2)|x-4|=1\Leftrightarrow\hoac{&(x-2)(x-4)=1\\&(x-2)(x-4)=-1.}
		\end{eqnarray*}		
		Khi đó ta suy ra nghiệm phương trình là $\hoac{&x=3-\sqrt{2} \,(\text{loại})\\&x=3+\sqrt{2}\\&x=3.}$\\
		Vậy tổng các nghiệm của phương trình là $S=3+3+\sqrt{2}=6+\sqrt{2}$ suy ra $a=6$, $b=1$.\\
		Do đó	giá trị của biểu thức $Q=a+3b=9$.	
	}
\end{ex}
\begin{ex}%[1D6H4-3]
	Bất phương trình $\log_2 \left(3x-1\right) > 3$ có nghiệm là
	\choice
	{$\dfrac{1}{3} < x < 3$}
	{\True $x > 3$}
	{$x < 3$}
	{$x > \dfrac{10}{3}$}
	\loigiai{
		$\log_2 \left(3x-1\right) > 3\Leftrightarrow \heva{&3x-1>0\\&3x-1>8}\Leftrightarrow x>3$.	}
\end{ex}
%%%==============EX_17============%%%
\begin{ex}%[1D6H4-2]
	Tìm tất cả các giá trị thực của $m$ để phương trình $5^x=m$ có nghiệm thực.
	\choice
	{$m\ge 1$}
	{$m\ge 0$}
	{\True $m > 0$}
	{$m\ne 0$}
	\loigiai{
		Vì $5^x>0$, $\forall x\in \mathbb{R}$ nên để phương trình $5^x=m$ có nghiệm thực thì $m>0$.	}
\end{ex}

\begin{ex}%[1D6H4-2]
	Nghiệm phương trình $\log_2 \left(3x-2\right)=4$ là
	\choice
	{\True $x=6$}
	{$x=7$}
	{$x=18$}
	{$x=4$}
	\loigiai{
		$\log_2 \left(3x-2\right)=4\Leftrightarrow \heva{&3x-2>0\\&3x-2=16}\Leftrightarrow x=6$.\\
		Vậy nghiệm phương trình là $x=6$.	}
\end{ex}
\begin{ex}%[1D6H4-2]
	Cho phương trình $3^{2x+1}-4\cdot3^x+1=0$ có 2 nghiệm $x_1$, $x_2$ trong đó $x_1 < x_2$. Chọn phát biểu đúng?
	\choice
	{$x_1+x_2=-2$}
	{\True $x_1+2x_2=-1$}
	{$x_1\cdot x_2=-1$}
	{$2x_1+x_2=0$}
	\loigiai{Đặt $t=3^x$, $t>0$. Khi đó phương trình đã cho được viết lại
		\[3t^2-4t+1=0\Leftrightarrow\hoac{&t=1\\&t=\dfrac{1}{3}}\Leftrightarrow \heva{&3^x=1\\&3^x=\dfrac{1}{3}}\Leftrightarrow\hoac{&x=0\\&x=-1.}\]
		Khi đó $x_1=-1$, $x_2=0$.		
	}
\end{ex}
\begin{ex}%[1D6H4-3]
	Gọi $S$ là tập hợp các nghiệm nguyên của bất phương trình $\left(\dfrac{1}{3} \right)^{\sqrt{x^2-3x-10}} > 3^{2-x}$.
	Tìm số phần tử của $S$.
	\choice
	{$9$}
	{$0$}
	{\True $8$}
	{$1$}
	\loigiai{
		\begin{eqnarray*}
			& & \left(\dfrac{1}{3} \right)^{\sqrt{x^2-3x-10}} > 3^{2-x}\\
			&\Leftrightarrow & \sqrt{x^2-3x-10}<x-2\\
			&\Leftrightarrow & \hoac{&x-2\ge0\\&x^2-3x-10>0\\&x^2-3x-10<(x-2)^2}\\
			&\Leftrightarrow & \hoac{&x\ge2\\&x<-2; x>5\\&x<14}\Leftrightarrow 5<x<14.
		\end{eqnarray*}	
		Vì $x$ nguyên nên $x\in \{6;7;8;9;10;11;12;13\}=S$. Do đó $S$ có số phần tử là $8$.}
\end{ex}

\begin{ex}%[1D6H4-3]
	Bất phương trình $\log _3\left(x^2-x+7\right)<2$ có tập nghiệm là khoảng $(a;b)$. Tính $b-a$.
	\choice
	{$b-a=-1$}
	{$b-a=-3$}
	{\True $b-a=3$}
	{$b-a=1$}
	\loigiai{
		Ta có $x^2-x+7=\left(x-\dfrac{1}{2}\right)^2+\dfrac{27}{4}>0,\ \forall x\in\mathbb{R}$.\\
		Khi đó $\log _3\left(x^2-x+7\right)<2 \Leftrightarrow x^2-x+7<9 \Leftrightarrow x^2-x-2<0 \Leftrightarrow-1<x<2$.\\
		Tập nghiệm là $S=(-1;  2)$. Suy ra $a=-1$, $b=2$.\\
		Vậy $b-a=2+1=3$.
	}
\end{ex}

\begin{ex}%[1D6V4-4]
	Tính tích các nghiệm của phương trình $\left(3+2\sqrt{2} \right)^{x^2-x+2}=\left(3-2\sqrt{2} \right)^{x^3-2}$
	\choice
	{\True $t=0$}
	{$t=2$}
	{$t=-1$}
	{$t=1$}
	\loigiai{
		\begin{eqnarray*}
			& & \left(3+2\sqrt{2} \right)^{x^2-x+2}=\left(3-2\sqrt{2} \right)^{x^3-2}\\
			&\Leftrightarrow & \left(3+2\sqrt{2} \right)^{x^2-x+2}=\dfrac{1}{\left(3+2\sqrt{2} \right)^{x^3-2}}\\
			&\Leftrightarrow & \left(3+2\sqrt{2} \right)^{x^2-x+2} \left(3+2\sqrt{2} \right)^{x^3-2}=1\\
			&\Leftrightarrow &\left(3+2\sqrt{2} \right)^{x^2-x+2+x^3-2}=1\\
			&\Leftrightarrow &	x^3+x^2-x=0\Leftrightarrow \hoac{&x=0 \\&x=\dfrac{-1\pm \sqrt{5}}{2}} \Rightarrow t=0.
		\end{eqnarray*}
	}
\end{ex}
\begin{ex}%[1D6V4-4]
	Tổng tất cả các nghiệm của phương trình $ 2^{x^2-2x-1}\cdot 3^{x^2-2x}=18 $ bằng
	\choice
	{$ 1 $}
	{$ -2 $}
	{\True $ 2 $}
	{$ -1 $}
	\loigiai{
		Ta có 
		\begin{eqnarray*}
			& & 2^{x^2-2x-1}\cdot 3^{x^2-2x}=18\\
			&\Leftrightarrow & 2^{x^2-2x-1}\cdot 3^{x^2-2x}=18\\
			&\Leftrightarrow & 2^{x^2-2x}\cdot\dfrac{1}{2}\cdot 3^{x^2-2x}=18\\
			&\Leftrightarrow & 6^{x^2-2x}=36\Leftrightarrow x^2-2x=2\Leftrightarrow x=1\pm \sqrt{3}.
		\end{eqnarray*}		
		Tổng các nghiệm của phương trình là $ \left(1+\sqrt{3}\right)+\left(1-\sqrt{3}\right)=2 $.
	}
\end{ex}

\begin{ex}%[1D6C4-5]
	Tìm tất cả giá trị của tham số $m$ để bất phương trình $\log \left(2x^2+3\right) > \log \left(x^2+mx+1\right)$ có tập nghiệm là $\mathbb{R}$.
	\choice
	{\True $-2< m < 2$}
	{$m < 2\sqrt{2}$}
	{$-2\sqrt{2} < m < 2\sqrt{2}$}
	{$m < 2$}
	\loigiai{
		Ta có $\log \left(2x^2+3\right) > \log \left(x^2+mx+1\right)$
		\[\Leftrightarrow \left\{\begin{array}{l} {x^2+mx+1 > 0} \\ {2x^2+3 > x^2+mx+1} \end{array}\right. \Leftrightarrow \left\{\begin{array}{l} {x^2+mx+1 > 0} \\ {x^2-mx+2 > 0} \end{array}\right. \left(*\right)
		\]
		Để bất phương trình $\log \left(2x^2+3\right) > \log \left(x^2+mx+1\right)$ có tập nghiệm là $\mathbb{R}$ thì hệ $\left(*\right)$ có tập nghiệm là $\mathbb{R}$
		\[\Leftrightarrow \left\{\begin{array}{l} {\Delta_1=m^2-4 < 0} \\ {\Delta_2=m^2-8 < 0} \end{array}\right. \Leftrightarrow-2< m < 2.
		\]
	}
\end{ex}
%%%==============EX_43============%%%
\begin{ex}%[1D6V4-5]
	Tìm tất cả các giá trị của tham số $m$ để bất phương trình $4^{x-1}-m\left(2^x+1\right) > 0$  nghiệm đúng với mọi $x\in \mathbb{R}$.
	\choice
	{$m\in \left(-\infty; 0\right)\cup \left(1;+\infty \right)$}
	{\True $m\in \left(-\infty; 0\right]$}
	{$m\in \left(0;+\infty \right)$}
	{$m\in \left(0; 1\right)$}
	\loigiai{
		Bất phương trình $4^{x-1}-m\left(2^x+1\right) > 0$. \hfill{(1)}\\		
		Đặt $t=2^x$, $t > 0$.\\		
		Bất phương trình (1) trở thành: $\dfrac{1}{4} t^2-m\left(t+1\right) > 0$ $\Leftrightarrow$ $t^2-4mt-4m > 0$. \hfill{(2)}\\		
		Đặt $f(t)=t^2-4mt-4m$.\\
		Đồ thị hàm số $y=f(t)$ có đồ thị là một Parabol với hệ số $a$ dương, đỉnh $I\left(2m;-4m^2-4m\right)$.\\
		Bất phương trình $(1)$ nghiệm đúng với mọi $x\in \mathbb{R}$ $\Leftrightarrow$ Bất phương trình $(2)$ nghiệm đúng với mọi $t > 0$ hay $f(t) > 0, \forall t > 0$.
		\begin{enumerate}[TH 1.]
			\item  $m\le 0$ $\Rightarrow$ $f(0)=-4m\ge 0$ $\Rightarrow$ $m\le 0$ thỏa mãn.
			\item  $m > 0$ $\Rightarrow$ $-4m^2-4m < 0$ nên $m > 0$ không thỏa mãn. \end{enumerate}		
		Vậy $m\le 0$.
	}
\end{ex}






\Closesolutionfile{ans}
\indapan{6}{ans/ans-1-C6-B4}
\cauds
\Opensolutionfile{ans}[ans/ans-1-C6-B4-DS]


\begin{ex}%[1D6V4-4]
	Cho phương trình $\log (x-1)^2=\log (x+1)$. Khi đó
	\choiceTF{Điều kiện $x > 1$}
	{Phương trình đã cho có chung tập nghiệm với phương trình $x^2-3x+\dfrac{9}{4}=0$}
	{\True Tổng các nghiệm của phương trình bằng $3$}
	{\True Biết phương trình có hai nghiệm $x_1, x_2 \left(x_1 < x_2 \right)$. Khi đó 3 số $x_1;x_2;6$ tạo thành một cấp số cộng}
	\loigiai{
		Điều kiện $\heva{&{(x-1)^2> 0} \\
			&{x+1> 0}.}$(*)
		\[\log (x-1)^2=\log (x+1)\Rightarrow (x-1)^2=x+1\Leftrightarrow x^2-3x=0\Leftrightarrow \hoac{&{x=0} \\
			&{x=3}}
		\]
		Thay lần lượt hai giá trị này vào $(*)$, ta thấy cả hai giá trị đều thoả mãn. Vậy phương trình có tập nghiệm là $S=\{0;3\}$.
		\begin{itemchoice}
			\itemch Sai
			\itemch Sai. Vì $x=0$ không có nghiệm $x^2-3x+\dfrac{9}{4}=0$.
			\itemch Đúng. Vì Tổng các nghiệm của phương trình bằng $0+3=3$.
			\itemch Đúng. Vì  $0;3;6$ tạo thành một cấp số cộng
		\end{itemchoice}	
	}
\end{ex}

\begin{ex}%[1D6H4-5]
	Giải các bất phương trình sau. Khi đó
	\choiceTF{$\log_2 (-x+3)\ge 1$ có nghiệm lớn nhất bằng $1$}
	{$\log _{\dfrac{1}{3}} (2x-2)\le 3$ có nghiệm bé nhất bằng $\dfrac{55}{54}$}
	{$\log_2 \left(x^2+5x+4\right) < 2$ có điều kiện nghiệm là $-4 < x <-1$}
	{$\log _{\dfrac{1}{9}} (-2x-1) > \log _{\dfrac{1}{9}} (x+1)$ tập nghiệm của bất phương này là: $S=\left(-\dfrac{2}{3};-\dfrac{1}{2} \right)$}
	\loigiai{
		\begin{itemchoice}
			\itemch 	Điều kiện $-x+3 > 0\Leftrightarrow x < 3$. $(*)$ Khi đó, do cơ số $2 > 1$ nên bất phương trình đã cho trở thành
			\[-x+3\ge 2^1 \Leftrightarrow x\le 1.\;
			\]
			Kết hợp với điều kiện $(*)$, ta được nghiệm của bất phương trình là $x\le 1$.
			\itemch Điều kiện $2x-2 > 0\Leftrightarrow x > 1$. $(*)$ Khi đó, do cơ số $0 < \dfrac{1}{3} < 1$ nên bất phương trình đã cho trở thành
			\[2x-2\ge \left(\dfrac{1}{3} \right)^3 \Leftrightarrow 2x\ge \dfrac{55}{27} \Leftrightarrow x\ge \dfrac{55}{54}.
			\]
			Kết hợp với điều kiện $(*)$, ta được nghiệm của bất phương trình là $x\ge \dfrac{55}{54}$
			\itemch Điều kiện $x^2+5x+4> 0\Leftrightarrow \hoac{&{x >-1} \\
				&{x <-4}.}$\\		
			Khi đó, do cơ số $2 > 1$ nên bất phương trình đã cho trở thành
			\[x^2+5x+4 < 2^2 \Leftrightarrow x^2+5x < 0\Leftrightarrow-5 < x < 0.
			\]
			Kết hợp với điều kiện $(*)$, ta được tập nghiệm của bất phương trình là:
			\[S=(-5;-4)\cup (-1;0).\;
			\]
			\itemch Điều kiện $\heva{&{-2x-1> 0} \\
				&{x+1> 0}}\Leftrightarrow \heva{&{x <-\dfrac{1}{2}} \\
				&{x >-1}}\Leftrightarrow-1< x <-\dfrac{1}{2}$. $(*)$\\
			Khi đó, do cơ số $0 < \dfrac{1}{9} < 1$ nên bất phương trình đã cho trở thành\\	
			$-2x-1 < x+1\Leftrightarrow x >-\dfrac{2}{3}$. Kết hợp điều kiện $(*)$, nghiệm của bất phương trình là $-\dfrac{2}{3} < x <-\dfrac{1}{2}$.
	\end{itemchoice}													}
\end{ex}
%%%==============EX_9============%%%
\begin{ex}%[1D6H4-5]
	Cho bất phương trình $\log _{0, 5} (x+1)^2 \le \log _{0, 5} 2x$, có tập nghiệm là $S=\left(a;b\right)$. Khi đó:
	\choiceTF{$a=0$}
	{$\left(a;b\right)\cap \left(3;2024\right)=\left(3;2024\right)$}
	{$A\left(a;0\right)$ là tọa độ đỉnh của parabol $\left(P\right):y=x^2+2$}
	{${\mathop{\lim}\limits_{x\to b}} \left(\dfrac{1}{x^3}+\dfrac{1}{x^2}+\dfrac{1}{x} \right)=3$}
	\loigiai{
		Điều kiện $\heva{&{(x+1)^2> 0} \\
			&{2x > 0}
			&} \Leftrightarrow \heva{&{x\ne-1} \\
			&{x > 0}} \Leftrightarrow x > 0.(*)$\\
		Khi đó, do cơ số $0 < 0, 5 < 1\;$ nên bất phương trình đã cho trở thành:
		\[(x+1)^2 \ge 2x\Leftrightarrow x^2+1\ge 0\Leftrightarrow x\in \mathbb{R}.
		\]
		Kết hợp với điều kiện $(*)$, ta được nghiệm của bất phương trình là $x > 0$.
	}
\end{ex}

\begin{ex}%[1D6V4-5]
	Cho bất phương trình $\left(\dfrac{1}{6} \right)^{x+2} \le \left(\dfrac{1}{36} \right)^{-x}$, có tập nghiệm là $S=\left[a;b\right)$. Khi đó:
	\choiceTF{Bất phương trình có chung tập nghiệm với $6^{-x-2} \le 6^{-2x}$}
	{${\mathop{\lim}\limits_{x\to b}} \left(3x^2+2\right)=b$}
	{$\left[a;b\right)\setminus \left(3;+\infty \right)=\left[-\dfrac{2}{3};3\right]$}
	{${\mathop{\lim}\limits_{x\to a}} \left(3x^2+2\right)=\dfrac{10}{3}$}
	\loigiai{
		$$\left(\dfrac{1}{6} \right)^{x+2} \le \left(\dfrac{1}{36} \right)^{-x} \Leftrightarrow 6^{-x-2} \le 6^{2x} \Leftrightarrow-x-2\le 2x\Leftrightarrow x\ge-\dfrac{2}{3} (\text{vì}\, 6 > 1).$$
		Một cách giải khác
		$$\left(\dfrac{1}{6} \right)^{x+2} \le \left(\dfrac{1}{36} \right)^{-x} \Leftrightarrow \left(\dfrac{1}{6} \right)^{x+2} \le \left(\dfrac{1}{6} \right)^{-2x} \Leftrightarrow x+2\ge-2x\Leftrightarrow x\ge-\dfrac{2}{3} (\text{vì}\,0 < \dfrac{1}{6} < 1).$$
		Vậy nghiệm của bất phương trình là $x\ge-\dfrac{2}{3}$.
	}
\end{ex}



\Closesolutionfile{ans}
\indapan{3}{ans/ans-1-C6-B4-DS}

\Opensolutionfile{ans}[ans/ans-1-C6-B4-KQ]
\caukq



\begin{ex}%[1D6V4-6]
	Anh Hưng gửi tiết kiệm khoản tiền 700 triệu đồng vào một ngân hàng với lãi suất $7\%$ / năm theo hình thức lãi kép kì hạn 12 tháng. Tính thời gian tối thiểu gửi tiết kiệm để anh Hưng thu được ít nhất 1 tỉ đồng (cả vốn lẫn lãi). Cho biết công thức lãi kép là $T=A\cdot (1+r)^n$, trong đó $A$ là tiền vốn, $T$ là tiền vốn và lãi nhận được sau $n$ năm, $r$ là lãi suất/năm.
	\shortans{$6$}	
	\loigiai{
		Ta có: $T\ge 1000\Leftrightarrow 700(1+7\%)^n \ge 1000\Leftrightarrow 1,07^n \ge \dfrac{10}{7}$
		\[\Leftrightarrow n\ge \log _{1,07} \left(\dfrac{10}{7} \right)\approx 5,27 \; (do\; 1,07 > 1).\;
		\]
		Vậy thời gian gửi tiết kiệm phải ít nhất 6 năm thì anh Hưng mới thu được ít nhât 1 tỉ đồng.
	}
\end{ex}

\begin{ex}%[1D6V4-4]
	Tổng tất cả các nghiệm của phương trình $16^{\tfrac{x+10}{x-10}}=0,125\cdot 8^{\tfrac{x+5}{x-15}}$ là
	\shortans{$20$}			
	\loigiai{
		Điều kiện xác định $\heva{&{x\ne 10} \\
			&{x\ne 15.}
		}$\\
		Phương trình đã cho viết lại $$2^{\tfrac{4x+40}{x-10}}=2^{-3} \cdot 2^{\tfrac{3x+15}{x-15}} \Leftrightarrow \dfrac{4x+40}{x-10}=\dfrac{60}{x-15} \Leftrightarrow 4x^2-80x=0
		\Leftrightarrow \hoac{&{x=0} \\
			&{x=20}.}$$
		Đối chiếu điều kiện, phương trình có nghiệm $x=0$ hoặc $x=20$.
	}
\end{ex}

\begin{ex}%[1D6H4-4]
	Tìm nghiệm của phương trình $\log_3 x+\log_9 x+\log _{27} x=11$.
	\shortans{$729$}		
	\loigiai{
		Với điều kiện $x > 0$, phương trình trở thành:
		\[\log_3 x+\dfrac{1}{2} \log_3 x+\dfrac{1}{3} \log_3 x=11\Leftrightarrow \log_3 x=6\Leftrightarrow x=3^6=729 > 0.
		\]
		Vậy nghiệm của phương trình là $x=729$.
	}
\end{ex}

\begin{ex}%[1D6V4-5]
	Số giá trị nguyên $x\in (-10;10)$ thỏa mãn bất phương trình $\log_3^2 (-x)-2\log _{\sqrt{3}} (-x)-2\log _{\tfrac{1}{3}} (-x)+1> 0$ là
	\shortans{$8$}
	\loigiai{
		Điều kiện: $-x > 0\Leftrightarrow x < 0.$\hfill{(*)}\\
		Khi đó bất phương trình tương đương
		\begin{eqnarray*}
			& & \log_3^2 (-x)-4\log _3 (-x)+2\log _3 (-x)+1> 0\\
			&\Leftrightarrow & \log_3^2 (-x)-2\log _3 (-x)+1> 0\\
			&\Leftrightarrow &\left( \log_3 (-x)-1\right)^2>0\\
			&\Leftrightarrow &\log_3 (-x)-1\not=0\Leftrightarrow x\not=-3.
		\end{eqnarray*}				
		Kết hợp với điều kiện $(*)$, ta được tập nghiệm của bất phương trình là
		\[S=(-\infty;-3)\cup (-3;0).\]
		Vì $x$ nguyên và $x\in (-10;10)$  nên $x\in \{-9;-8;-7;-6;-5;-4;-2;-1\}$.\\ Do đó có $8$ giá trị nguyên của $x$ thỏa bài toán.
	}
\end{ex}


\begin{ex}%[1D6C4-5]
	Cho bất phương trình $ \left(3^{x^2 - x} - 9 \right )\left ( 2^{x^2} - m \right ) \leq 0 $. Tìm số giá trị nguyên của $ m $ để bất phương trình đã cho có đúng $ 5 $ nghiệm nguyên.
	\shortans{$65024$}	
	\loigiai{
		Ta thấy $ 3^{x^2-x} - 9 < 0 \Leftrightarrow x^2 - x - 2 < 0 \Leftrightarrow -1< x < 2  $.\\
		Do vậy, bài toán cần tìm $ m $ nguyênđể $ \heva{& x \leq - 1 \lor x \geq 2 \\ & m \geq 2^{x^2}} \quad (*) $ có đúng $ 5 $ nghiệm $ x $ nguyên.\\
		Ta có bảng biến thiên của hàm số $ f(x) = 2^{x^2} $
		\begin{center}
			\begin{tikzpicture}[scale=0.8, font=\footnotesize, line join=round, line cap=round,>=stealth,yscale=.8,xscale=1.3]
				\def\cot{7} % số nhãn chiều dài
				\def\hang{4} % số nhãn chiều cao
				\draw[shift={(-.5,.5)}] 
				(0,0) rectangle +(\cot+1,-\hang-1)
				(0,-1)--+(0:\cot+1) 	
				(1,0)--+(-90:\hang+1);
				\path 
				(0,0) node{$x$} 
				(1,0) node{$-\infty$}
				(3,0) node{$-1$}
				(5,0) node{$2$}
				(7,0) node{$+\infty$}										 	
				(0,-2) node{$f(x)$}
				(1,-1) node (dvcL) {$+\infty$}												
				(7,-1) node (dvcR) {$+\infty$}
				(2.8,-4) node (CT) {$2$}
				(5.2,-3) node (CT1) {$2^4$}
				;
				\draw[->] (dvcL)--(CT);	
				\draw[->] (CT1)--(dvcR);
				\fill[pattern=north east lines,pattern color=gray] (3,-.5) rectangle (5,-4.5);
				\draw (3,-.5)--(3,-4.5) (5,-.5)--(5,-4.5);
			\end{tikzpicture}
		\end{center}		
		Từ bảng biến thiên, ta thấy $ (*) $ có đúng $ 5 $ nghiệm nguyên là $ \{-3;-2;-1;2;3 \} $, khi đó ta được $ 2^9 \leq m < 2^{16} $.\\
		Vậy có $ 2^{16} - 2^9 = 65024 $ giá trị nguyên của tham số $ m $ thoả đề bài.
	}
\end{ex}

\begin{ex}%[1D6C4-5]
	Cho $x$, $y$ là các số thực dương thỏa mãn $\ln x+\ln y\ge \ln (x^2+y)$. Tìm giá trị nhỏ nhất của $P=x+y$ (làm tròn đến hàng phần trăm).
	\shortans{$5{,}82$}
	\loigiai{
		Ta có $\ln x+\ln y\ge \ln (x^2+y)\Leftrightarrow xy\ge x^2+y\Leftrightarrow y(x-1)\ge x^2>0$.\\
		Suy ra $x>1$ và $y>0$.\\
		Khi đó $y\ge \dfrac{x^2}{x-1}$.\\
		Suy ra $P=x+y\ge x+\dfrac{x^2}{x-1}=2(x-1)+\dfrac{1}{x-1}+3\ge 2\sqrt{2(x-1)\cdot \dfrac{1}{x-1}}+3=2\sqrt{2}+3$.\\
		Vậy giá trị nhỏ nhất của biểu thức $P$ là $3+2\sqrt{2}\approx5{,}82$.
	}
\end{ex}

\Closesolutionfile{ans}
\indapan{6}{ans/ans-1-C6-B4-KQ}